\chapter{면 구조와 꼭짓점 나누기}
\label{chap:faces}

\section{면 구조 (조합적 조르당)}

이 절의 도구는 ``종수 0의 경직성''이다 (\ref{chap:routes}장). 연결 성분 수는 그대로 두면서
\(\chi\)를 2 늘리는 수술은 종수 0 회전 체계에서 불가능하다 (정리~\ref{thm:euler}).

\begin{lemma}[같은 면의 두 모서리에서 꼭짓점 쪼개기]
\label{lem:split2}
\uses{def:rotation, def:eulerChar}
\lean{Kuratowski.splitGraph, Kuratowski.splitDartEquiv, Kuratowski.RotationSystem.split,
  Kuratowski.RotationSystem.split_facePerm, Kuratowski.RotationSystem.eulerChar_split,
  Kuratowski.SplitLinked, Kuratowski.card_connectedComponent_le_split,
  Kuratowski.card_connectedComponent_split_le, Kuratowski.card_connectedComponent_split_of_linked,
  Kuratowski.card_connectedComponent_split_of_not_linked,
  Kuratowski.RotationSystem.not_splitLinked_of_eulerChar}
\leanok
유한하고 같음을 판정할 수 있는 꼭짓점 타입 \(V\) 위의 그래프 \(G\)와 회전 체계 \(R\)(회전 \(\rho\), 다트 뒤집기 \(\sigma\),
면 순열 \(\varphi = \rho\sigma\))에서, 같은 꼭짓점 \(w\)로 들어오는 서로 다른 두 다트 \(e_1 \ne e_2\)
(\(e_1, e_2\)의 끝점이 \(w\))가 \(\varphi\)의 같은 순환에 있다고 하자. 곧 면이 두 모서리
\((e_i, \varphi(e_i))\)에서 \(w\)를 지난다. 그러면 \(w\)를 둘로 쪼갠 그래프 \(G'\)(꼭짓점 타입
\(\mathrm{Option}\,V\), 새 꼭짓점은 \(\mathrm{none}\))와 회전 체계 \(R'\)이 있어
\[
  \chi(R') = \chi(R) + 2
\]
이다. 또 \(G'\)의 연결 성분 수 \(c'\)는 \(c \le c' \le c + 1\)이다. 아래의 두 순환 \(A\), \(B\)에 대해
어떤 \(a \in A\), \(b \in B\)가 있어 \(G - w\)에서 \(a\)의 끝점과 \(b\)의 끝점이 이어지면
\(c' = c\)이고, 그런 쌍이 하나도 없으면 \(c' = c + 1\)이다. 종수 0은 가정하지 않는다.
\bfcheck{\texttt{[split2]}: 꼭짓점 5개 이하의 모든 그래프와 모든 회전 체계에서, 같은 면에 있고 같은
꼭짓점으로 들어오는 서로 다른 두 다트 291,135쌍. 위 구성으로 쪼갠 회전이 꼭짓점마다 한 순환이고
\(n' = n + 1\), \(E' = E\), \(I' = I\), \(O' = O + 1\), \(\chi' = \chi + 2\)이며 성분 조건이 맞는다.
\texttt{[split2-음성]}: 두 다트가 다른 면에 있는 328,839쌍에서는 모두 \(O' = O - 1\)이라 같은 면
가정이 필요하다. \texttt{[split2-절단점]}: 경로 \(0\)--\(1\)--\(2\)의 가운데를 쪼개면 \(c' = c + 1\)이다.
유한 검사이며 일반 증명을 대신하지 않는다.}
\end{lemma}

\begin{proof}
\uses{lem:swap, thm:euler}
\leanok
\(\bar e_1, \bar e_2\)는 \(w\)에서 나가는 다트이고 \(\bar e_1 \ne \bar e_2\)다. 회전을
\[
  \rho' = \rho\,(\bar e_1\ \bar e_2)
\]
로 바꾼다. 두 다트가 \(w\)의 한 \(\rho\)-순환에 있으므로 \(w\)의 다트는 \(\rho'\)에서 정확히 두 순환
\(A \ni \bar e_1\), \(B \ni \bar e_2\)로 나뉘고, 둘 다 비어 있지 않다 (보조정리~\ref{lem:swap}).
\(w\)가 아닌 꼭짓점의 다트에서는 \(\rho' = \rho\)다.
\begin{itemize}
  \item \textbf{새 그래프.} 다트 \(d\)를
        \[
          d \mapsto \bigl(\,[d \in B]\ ?\ \mathrm{none} : d_{\mathrm{fst}},\ \
          [\bar d \in B]\ ?\ \mathrm{none} : d_{\mathrm{snd}}\,\bigr)
        \]
        로 옮기고(\(\mathrm{Option}\,V\)에서 \(V\)의 점은 \(\mathrm{some}\)으로 넣는다), 그 상들을 변으로 하는
        그래프를 \(G'\)라 하자. \(d\)와 \(\bar d\)가 둘 다 \(B\)에 있으면 \(d\)의 두 끝이 모두 \(w\)이므로
        고리는 생기지 않는다. 변 \(wu\)는 다트 \((w,u)\)가 \(A\)에 있으면 \(\mathrm{some}\,w\)--\(\mathrm{some}\,u\),
        \(B\)에 있으면 \(\mathrm{none}\)--\(\mathrm{some}\,u\)가 되고 두 경우는 겹치지 않으므로 \(G'\)도 단순
        그래프다. \(\mathrm{none}\)과 \(\mathrm{some}\,w\) 사이의 변은 없다.
  \item \textbf{다트 대응.} 위 사상은 \(G\)의 다트와 \(G'\)의 다트의 일대일 대응이고 뒤집기와 교환한다.
        \(R'\)의 회전은 이 대응으로 옮긴 \(\rho'\)다. \(B\)는 \(\rho'\)-순환이므로 시작점이 보존되고,
        \(\mathrm{none}\)의 다트는 \(B\), \(\mathrm{some}\,w\)의 다트는 \(A\)라는 한 순환을 이룬다.
  \item \textbf{면.} \(\varphi' = \rho'\sigma = \rho\sigma\,\sigma(\bar e_1\ \bar e_2)\sigma
        = \varphi\,(e_1\ e_2)\)다. \(e_1, e_2\)가 같은 \(\varphi\)-순환에 있으므로 그 순환이 갈라져
        \(O' = O + 1\)이다 (보조정리~\ref{lem:swap}).
  \item \textbf{세기.} \(n' = n + 1\)(\(\mathrm{Option}\,V\)), 다트 대응으로 \(E' = E\)다. \(w\)의 두 반쪽은
        각각 \(A\), \(B\)의 다트를 가지므로 고립점이 아니고 나머지 점은 그대로라 \(I' = I\)다. 따라서
        \(\chi' = (n+1) - E + (O+1) + I = \chi + 2\)다.
  \item \textbf{성분.} \(\mathrm{none}\)을 \(w\)로 보내는 사상은 \(G' \to G\) 준동형이므로 꼭짓점 쪼개기는
        성분을 합치지 않고, 갈라지는 성분은 \(w\)의 성분 하나뿐이라 \(c \le c' \le c + 1\)이다.
        \(\mathrm{some}\,w\)와 \(\mathrm{none}\)이 \(G'\)에서 이어지는 것은 어떤 \(a \in A\), \(b \in B\)에 대해
        \(a\)의 끝점과 \(b\)의 끝점이 \(G - w\)에서 이어지는 것과 같다. \(\mathrm{some}\,w\)에서 \(\mathrm{none}\)으로
        가는 길은 \(A\)의 다트로 나가 \(G - w\) 안을 지나 \(B\)의 다트로 들어오기 때문이다. 그런 쌍이 있으면
        \(c' = c\)이고, 없으면 \(c' = c + 1\)이다.
\end{itemize}
예: 경로 \(a\)--\(w\)--\(b\)에서 면 \((a,w) \to (w,b) \to (b,w) \to (w,a)\)는 \(w\)를 두 번 지난다.
\(e_1 = (a,w)\), \(e_2 = (b,w)\)로 쪼개면 \(G'\)는 \(a\)--\(w\)와 \(\mathrm{none}\)--\(b\)이고, \(n\)은
\(3 \to 4\), \(O\)는 \(1 \to 2\), \(\chi\)는 \(2 \to 4\)이며 \(w\)가 절단점이라 \(c\)는 \(1 \to 2\)다.

이 보조정리는 보조정리~\ref{lem:J1}에서 쓰인다. \(G\)가 2-연결이면 \(G - w\)에서 \(w\)의 이웃들이 모두
이어지므로 위 조건을 만족하는 쌍 \(a \in A\), \(b \in B\)가 있어 \(c' = c\)이고, \(R\)이 종수 0이면 \(\chi(R') = 2c + 2 > 2c'\)이 되어
정리~\ref{thm:euler}에 모순이다.

Lean의 \texttt{SplitLinked B w}는 \(w\)에서 시작하는 \(B\) 밖 다트와 \(B\) 안 다트의 두 끝점이
\(G-w\)에서 이어지는 존재 조건이다. \texttt{SplitVertex.lean}에서 그래프와 다트 대응, 회전·면 순열 식,
고립점·성분 수 조건까지 형식화했다.
\end{proof}


\begin{lemma}[2-연결 평면 그래프의 면 경계는 사이클]
\label{lem:J1}
\uses{def:rotation, def:eulerChar, def:conn3}
\lean{Kuratowski.RotationSystem.eq_of_sameCycle_facePerm_of_fst_eq}
\leanok
유한하고 같음을 판정할 수 있는 \(V\) 위에서 \(G\)가 2-연결이고 회전 체계 \(R\)이 종수 0(\(\chi(R) = 2c\))이면, 면 순열의 한 순환 위 다트들의
시작점은 서로 다르다.
\bfcheck{\texttt{[J1]}: 회전 체계 811개. \texttt{[J1-음성]}: 종수 0 조건을 빼면 \(K_4\)에서 반례가 나온다
(종수 0이 아닌 회전 체계 14개 모두).}
\end{lemma}

\begin{proof}
\uses{lem:split2, thm:euler}
\leanok
같은 면의 다트 \(d\ne e\)가 같은 시작점 \(w\)를 가진다고 하자. 들어오는 두 다트는
\(e_1=\varphi^{-1}(d)\), \(e_2=\varphi^{-1}(e)\)로 잡는다. 둘은 서로 다르고 같은 면에 있으며 끝점이
\(w\)다. 역다트 \(\bar d,\bar e\)를 고르면 같은 면이 보장되지 않으므로 면 순열의 역상을 쓴다.
2-연결성에서 \(G-w\)의 두 이웃이 이어져 보조정리~\ref{lem:split2}의 연결 조건이 성립한다.
따라서 \(\chi\)는 2 늘고 성분 수는 그대로라 오일러 부등식에 모순이다.

Lean 진술은 같은 면 순환에서 시작점이 같은 두 다트가 같다는 형태다. 실제 순환 수열은 다음
보조정리~\ref{lem:face_cycle}에서 구성한다. 지운 꼭짓점의 이웃 순서는
보조정리~\ref{lem:J2}에서 이 면 순환 계산에 연결한다.
\end{proof}

\begin{lemma}[면 경계의 순서를 보존하는 사이클]
\label{lem:face_cycle}
\uses{def:rotation, def:eulerChar, def:conn3, def:cycseq}
\lean{Kuratowski.RotationSystem.faceLength, Kuratowski.RotationSystem.three_le_faceLength,
  Kuratowski.RotationSystem.exists_cyclicSequence_face}
\leanok
유한하고 같음을 판정할 수 있는 \(V\) 위의 2-연결 그래프 \(G\), 종수 0 회전 체계 \(R\), 다트 \(d\)가
주어지면, 어떤 \(m\ge3\)과 \(C:\mathrm{CyclicSequence}\,G\,m\)이 있어 모든 자연수 \(k\)에 대해
\[
 C(k\bmod m)=\bigl(\varphi^k(d)\bigr)_{\mathrm{fst}}
\]
이다. 따라서 \(C(0)=d_{\mathrm{fst}}\),
\(C(1)=d_{\mathrm{snd}}\)이며 면 경계의 시작점과 정방향을 보존한다. 길이 3 이상의 하한에는 종수 0이
필요하지 않다.
\end{lemma}

\begin{proof}
\uses{lem:J1}
\leanok
면 순열의 최소 주기 \(m=\mathrm{faceLength}_R(d)\)로 길이를 고른다. 유한 순열의 최소 주기는 양수다. 길이가 1이면 다트의 두 끝이 같아 고리 금지에 어긋난다. 길이가 2이면
\(\varphi(d)=\bar d\)와 \(\rho(d)=d\)에서 \(d\)의 시작점은 이웃이 하나뿐이다. 그 이웃을 지우면 고립되므로
지지점이 세 개 이상인 2-연결성에 어긋난다. \(i\mapsto(\varphi^{i.\mathrm{val}}d)_{\mathrm{fst}}\)로 수열을
만들면 단사성은 보조정리~\ref{lem:J1}과 최소 주기 아래의 순열 반복 단사성에서 온다. 다음 점과의 인접은
\((\varphi e)_{\mathrm{fst}}=e_{\mathrm{snd}}\)에서 나온다.
\end{proof}

\begin{definition}[꼭짓점 지우기]
\label{def:delVertex}
\uses{def:rotation, thm:del}
\lean{Kuratowski.RotationSystem.delVertex, Kuratowski.RotationSystem.delVertex_rot_toProd,
  Kuratowski.RotationSystem.delVertex_rot_of_snd_ne,
  Kuratowski.RotationSystem.delVertex_rot_of_snd_eq,
  Kuratowski.RotationSystem.delVertex_facePerm_toProd, Kuratowski.RotationSystem.eq_delVertex,
  Kuratowski.RotationSystem.snd_rot_rot_ne, Kuratowski.RotationSystem.eulerChar_delVertex,
  Kuratowski.IsPlanar.deleteVertex}
\leanok
유한한 꼭짓점 타입 \(V\)와 \texttt{DecidableEq V}에서, \(v\)의 변을 모두 지운 그래프
\texttt{deleteVertices G \{v\}}의 회전 체계 \(R.\mathrm{delVertex}\,v\)를 만든다. \(v\) 자체는
고립점으로 남으며, 처음부터 고립점이어도 된다.

원래 회전을 \(\rho\)라 하자. 살아남은 다트 \(d\)는 양 끝점이 \(v\)가 아니다. 대응하는 원래 다트도
\(d\)로 쓰면, 새 회전은 \((\rho d).\mathrm{snd}=v\)일 때 \(\rho^2d\), 아니면 \(\rho d\)다.
단순 그래프라 각 \(u\ne v\)에서 지워지는 다트 \((u,v)\)는 하나뿐이므로 한 칸만 건너뛰면 된다.
실제로 \(d\)가 \(v\)로 가지 않고 \(\rho d\)가 \(v\)로 가면 \(\rho^2d\)는 \(v\)로 가지 않는다.
면 순열에도 역다트를 대입한 같은 공식이 성립한다. 이 공식은 새 회전 체계를 유일하게 정한다.

구성은 \(v\)의 이웃 유한 집합에 대한 귀납법으로 정리~\ref{thm:del}의 변 삭제를 반복한다.
각 단계에서 위 공식을 유지하며, \(R\)이 종수 0이면 이 구성된 회전 체계도 종수 0이다.
따라서 평면성도 유지되며, 증인은 다른 회전 체계를 고르는 것이 아니라 \(R.\mathrm{delVertex}\,v\)다.
보조정리~\ref{lem:J2}는 이 실제 회전 체계의 특정 면에서 이웃의 순환 순서를 얻는다.
\end{definition}

\begin{lemma}[지운 꼭짓점 둘레의 면]
\label{lem:J2}
\uses{def:rotation, def:eulerChar, def:conn3, def:delVertex, def:cycseq}
\lean{Kuratowski.RotationSystem.exists_cyclicSequence_delVertex_neighbors,
  Kuratowski.RotationSystem.minimalPeriod_rot_inv, Kuratowski.RotationSystem.rot_inv_pow_ncard,
  Kuratowski.RotationSystem.exists_rot_inv_pow_snd, Kuratowski.RotationSystem.sector_iterate,
  Kuratowski.RotationSystem.sector_iterate_next,
  Kuratowski.RotationSystem.faceLength_delVertex_eq_sectorStart,
  Kuratowski.RotationSystem.eq_faceLength_of_cyclicSequence}
\leanok
유한하고 같음을 판정할 수 있는 꼭짓점 타입 \(V\) 위의 그래프 \(H\)와 종수 0 회전 체계 \(R\)(회전 \(\rho\),
면 순열 \(\varphi\))을 생각한다. 꼭짓점 \(v\)의 이웃 수를 \(N = |N_H(v)|\)라 하고 \(N \ge 2\)라 하자.
지지 집합에서 \(v\)를 뺀 \(S_H \setminus \{v\}\)가 \(H - v\)(\texttt{deleteVertices H \{v\}})에서
2-연결이라 하자 (정의~\ref{def:conn3}). \(v\)에서 나가는 다트 \(d_0\)를 하나 고르고
\[
  d_j = \rho^{-j} d_0, \qquad u_j = (d_j)\text{의 끝점} \qquad (0 \le j < N)
\]
이라 두자. 곧 \(u_0, u_1, \ldots, u_{N-1}\)은 \(v\) 둘레의 회전을 거꾸로 돈 이웃 열이다.
\(\varphi(d_0)\)와 양 끝이 같은 \(H - v\)의 다트 \(e\)가 존재한다 (\(\varphi(d_0)\)는 \(u_0\)에서 출발해 \(v\)로 가지 않는다).
그러면 길이 \(m\)의 순환 꼭짓점 수열 \(C\)(정의~\ref{def:cycseq})가 있어
\[
  C(k\bmod m) = \bigl((R.\mathrm{delVertex}\,v).\varphi^{k}\, e\bigr)\text{의 시작점} \qquad (k \in \mathbb{N})
\]
이다. 곧 \(C\)는 실제 회전 체계 \(R.\mathrm{delVertex}\,v\)(정의~\ref{def:delVertex})에서 \(e\)가 속한 면이며
새 회전 증인을 고르지 않는다. 또 자연수 위치 \(p_0 = 0 < p_1 < \cdots < p_{N-1} < m\)이 있어
\(C(p_j) = u_j\)이다. 곧 \(e\)에서 출발해 그 면을 따라가면 이웃을 \(v\) 둘레 회전의 역순으로 정확히 한 번씩 만난다.

이 진술은 ``\(H - v\)의 어떤 면이 \(u_0, \ldots, u_{N-1}\)을 이 순환 순서대로 지나는 사이클이다''를 강화한 것이다.
그 존재 진술도 참이며 \texttt{[J2]}가 검사한다. 다만 \(H - v\)가 사이클이면 두 면이 모두 이웃을 지나고
반대쪽 면은 정순으로 지나므로, 존재 진술은 방향을 바꾼 음성 대조를 구별하지 못한다. 그래서 면을 \(v\)가
있던 자리의 면, 곧 \(e\)의 면으로 고정한다. 증명이 실제로 만드는 면도 이 면이다.

가정을 ``\(H - v\)가 2-연결''(곧 \(H - v\)의 지지 집합 기준)로 약하게 하면 거짓이다. 변이
\(\{01, 12, 20, 30, 31, 34\}\)이고 \(v = 3\)인 평면 그래프에서 \(H - v\)는 삼각형 \(012\)와 고립점 \(3\), \(4\)라서
2-연결이지만, 이웃 \(4\)가 \(H - v\)의 고립점이라 이웃 \(\{0, 1, 4\}\)를 지나는 사이클이 없다. 지금
가정에서는 \(\{0, 1, 2, 4\}\)가 \(H - v\)에서 이어져 있지 않으므로 이 경우가 빠진다.
\bfcheck{\texttt{[J2]}: 1,892개 경우(어떤 면이 사이클이고 이웃을 역순의 순환 순서로 지남).
\texttt{[J2-위치]}: 같은 경우의 모든 시작 이웃 4,620개에서, 면을 \(e = \varphi(v, s)\)로 고정하면 꼭짓점이 겹치지 않고
\(u_j\)의 위치가 \(0\)에서 엄격히 증가한다. \texttt{[J2-위치-음성]}: 같은 고정 면에 정순 \(\rho^{j} d_0\)의 끝점을
넣으면 \(N \ge 3\)인 2,284개 모두에서 실패한다. \texttt{[J2-구간]}: 같은 4,620개에서 아래 증명의 구간 분해,
곧 \(L_j \ge 3\), \(\varphi^{L_j-1} d_j = \bar d_{j+1}\), \(P_j\)가 \(v\)에 닿지 않음, \(d_j\)들의 면이 서로 다름,
\(m = \sum_j (L_j - 2)\), \(p_j = \sum_{t<j} (L_t - 2)\)가 성립한다. \texttt{[J2-반례]}: 위 그래프에서 약한 가정은
성립하고 결론은 틀리며(종수 0 회전 체계 4개), 지금 가정은 성립하지 않는다. 유한 검사이며 일반 증명을
대신하지 않는다.}
\end{lemma}

\begin{proof}
\uses{lem:J1, lem:face_cycle}
\leanok
\(\varphi'\)을 \(R.\mathrm{delVertex}\,v\)의 면 순열이라 하자. \(H - v\)의 다트와 \(v\)에 닿지 않는 \(H\)의 다트는
양 끝으로 대응하므로, 아래에서 두 다트를 양 끝이 같으면 같은 것으로 적는다.
\begin{itemize}
  \item \textbf{\(H - v\)와 \(H\)의 2-연결.} \(S_H \setminus \{v\}\)는 점이 3개 이상이고 \(H - v\)에서
        이어져 있으므로 각 점이 \(H - v\)에서 이웃을 갖는다. 따라서 \((H - v)\)의 지지 집합은
        \(S_H \setminus \{v\}\)이고 \(H - v\)는 2-연결이다 (\texttt{TwoConnected}).
        \(H\)도 2-연결이다. 아무것도 지우지 않으면 \(S_H \setminus \{v\}\)가 이어져 있고 \(v\)는 이웃으로 거기에
        붙는다. \(v\)를 지우면 가정 그대로다. 다른 점 \(z\)를 지우면 \(S_H \setminus \{v, z\}\)가 \(H - \{v, z\}\)에서
        이어져 있고, \(N \ge 2\)이므로 \(v\)는 \(z\)가 아닌 이웃을 통해 거기에 이어진다.
  \item \textbf{\(v\) 둘레의 회전.} \(v\)에서 나가는 다트는 \(\rho\)의 한 순환을 이루고(회전 체계의 조건)
        그 수는 \(N\)이다. 따라서 \(d_0, \ldots, d_{N-1}\)은 서로 다르고 \(d_N = d_0\)이다.
        \(\rho^{-1} d_j = d_{j+1}\)이다.
  \item \textbf{\(H\)의 면 조각.} \(L_j\)를 \(d_j\)의 면 길이(\(\varphi\)의 최소 주기)라 하자.
        \(H\)가 2-연결이므로 \(L_j \ge 3\)이다 (보조정리~\ref{lem:face_cycle}의 길이 하한).
        \(\varphi(x) = d_j\)이면 \(\rho(\bar x) = d_j\)라 \(x = \overline{\rho^{-1} d_j} = \bar d_{j+1}\)이므로
        \(\varphi^{L_j - 1} d_j = \bar d_{j+1} = (u_{j+1}, v)\)이다. \(H\)는 2-연결이고 \(R\)은 종수 0이므로
        보조정리~\ref{lem:J1}에서 이 면의 다트 \(\varphi^i d_j\) (\(0 \le i < L_j\))의 시작점은 서로 다르다.
        시작점이 \(v\)인 것은 \(i = 0\)뿐이고, \(\varphi^i d_j\)의 끝점은 \(\varphi^{i+1} d_j\)의 시작점이다.
        따라서 \(\ell_j = L_j - 2 \ge 1\)개의 다트
        \[
          P_j = \bigl(\varphi^{1} d_j, \varphi^{2} d_j, \ldots, \varphi^{\ell_j} d_j\bigr)
        \]
        는 모두 \(v\)에 닿지 않는 \(H - v\)의 다트이고, \(P_j\)는 \(u_j\)에서 \(u_{j+1}\)로 가는 걸음이다.
        특히 \(\varphi(d_0)\)는 \(v\)에 닿지 않으므로 \(e\)가 있다.
  \item \textbf{조각은 서로소.} \(j \ne j'\)이면 \(d_j \ne d_{j'}\)이고 둘의 시작점은 모두 \(v\)이므로,
        보조정리~\ref{lem:J1}에서 \(d_j\)와 \(d_{j'}\)은 \(\varphi\)의 서로 다른 순환에 있다. \(P_j\)의 다트는
        \(d_j\)의 면에 있으므로 서로 다른 \(j\)의 조각은 다트를 공유하지 않는다. 한 조각 안의 다트도
        최소 주기 아래의 반복이라 서로 다르다. 따라서 \(M = \ell_0 + \cdots + \ell_{N-1}\)개의 다트가 모두 다르다.
  \item \textbf{지운 회전의 계산.} \(H - v\)의 다트 \(x\)에 대해 \(\varphi(x)\)가 \(v\)로 가지 않으면
        \(\varphi'(x) = \varphi(x)\)이고, \(\varphi(x) = (w, v)\)이면 \(\varphi'(x) = \rho((w, v)) = \varphi(v, w)\)다
        (\texttt{delVertex\_facePerm\_toProd}). \(P_j\)의 마지막이 아닌 다트에서는 다음 다트가 \(P_j\) 안에 있어
        첫 경우이고, 마지막 다트 \(\varphi^{\ell_j} d_j\)에서는 \(\varphi\)가 \(\bar d_{j+1}\)이라 둘째 경우다.
        따라서 \(s_j = \ell_0 + \cdots + \ell_{j-1}\)이라 두면 \(0 \le i < \ell_j\)에서
        \[
          \varphi'^{\,s_j + i}\, e = \varphi^{1 + i} d_j
        \]
        이다 (\(j\)와 \(i\)에 대한 귀납, \(P_0\)의 첫 다트가 \(e\)). \(d_N = d_0\)이므로
        \(\varphi'^{\,M} e = \varphi(d_0) = e\)이다.
  \item \textbf{면 길이와 위치.} 위 두 항목에서 \(\varphi'^{\,k} e\) (\(0 \le k < M\))는 서로 다르고
        \(\varphi'^{\,M} e = e\)이므로 \(e\)의 \(\varphi'\) 최소 주기는 정확히 \(m = M = \sum_{j<N} (L_j - 2)\)다.
        \(R.\mathrm{delVertex}\,v\)는 종수 0이고(정의~\ref{def:delVertex}) \(H - v\)는 2-연결이므로
        보조정리~\ref{lem:face_cycle}을 \(e\)에 적용하면 길이 \(m\)의 순환 꼭짓점 수열 \(C\)가 있어
        \(C(k)\)는 \(\varphi'^{\,k} e\)의 시작점이다. \(p_j = s_j\)로 두면 \(p_0 = 0\),
        \(p_{j+1} = p_j + \ell_j > p_j\), \(p_{N-1} < p_{N-1} + \ell_{N-1} = m\)이고
        \(C(p_j) = (\varphi d_j)\text{의 시작점} = u_j\)이다.
\end{itemize}
Lean의 \texttt{DelVertexFace.lean}은 위 이웃 열거, 두 2-연결 유도, 조각 반복과 최소 주기 식을
형식화한다. 기존 \texttt{exists\_cyclicSequence\_face}의 인터페이스는 그대로 두고, 그 수열의 길이가
면의 최소 주기와 같다는 정리로 길이를 복원한다. 진술에 적은 수학적 가정 외에는
꼭짓점 타입의 유한성과 \texttt{DecidableEq V}만 쓴다.
\end{proof}


\begin{definition}[순환 꼭짓점 수열]
\label{def:cycseq}
\lean{Kuratowski.CyclicSequence, Kuratowski.CyclicSequence.rotate,
  Kuratowski.CyclicSequence.reverse, Kuratowski.CyclicSequence.neZero,
  Kuratowski.CyclicSequence.adj_last_zero,
  Kuratowski.CyclicSequence.triangle, Kuratowski.CyclicSequence.triangle_closing_edge,
  Kuratowski.CyclicSequence.not_nonempty_of_lt_three, Kuratowski.CyclicSequence.not_nonempty_bot}
\leanok
사이클은 단사 함수 \(c : \mathbb{Z}/m \to V\) (\(m \ge 3\), \(c(i) \sim c(i + 1)\))로 나타낸다. 사이클을
구간 단위로 다루기 위한 표현이다 (결정~D3). Lean의 \texttt{CyclicSequence G m}은 길이 조건,
\texttt{vertex : ZMod m} \(\hookrightarrow\) \texttt{V}, 이웃한 위치의 인접 조건을 묶는다.
모듈러 덧셈이므로 마지막 위치와 첫 위치 사이의 변도 포함한다. 유도 사이클일 필요는 없다.
\texttt{rotate}는 시작점을 옮기고 \texttt{reverse}는 방향을 뒤집는다. 마지막 위치 \(-1\)과 첫 위치
\(0\)의 인접은 \texttt{adj\_next}에서 나오며, 삼각형 예제는 \(2 + 1 = 0\)을 써서 닫는 변을 확인한다.
길이 \(0,1,2\)와 변이 없는 그래프는 배제된다.
\end{definition}

\begin{lemma}[순서를 보존하는 사이클과 두 호]
\label{lem:cycseq_walks}
\uses{def:cycseq, def:paths}
\lean{Kuratowski.CyclicSequence.toWalk, Kuratowski.CyclicSequence.toWalk_isCycle,
  Kuratowski.CyclicSequence.length_toWalk, Kuratowski.CyclicSequence.getVert_toWalk,
  Kuratowski.CyclicSequence.arc, Kuratowski.CyclicSequence.arc_isPath,
  Kuratowski.CyclicSequence.length_arc, Kuratowski.CyclicSequence.mem_support_arc_iff,
  Kuratowski.CyclicSequence.mem_walkInterior_arc, Kuratowski.CyclicSequence.disjoint_arc_interiors}
\leanok
순환 수열 \(c\)는 \(c(0)\)에서 출발해 돌아오는 길이 \(m\)의 Mathlib \texttt{IsCycle}을 준다.
\(0 \le i \le m\)에서 이 걸음의 \(i\)번째 점은 \(c(i)\)다. 위치 \(a\)에서 \(b\)까지의 정방향 호는
길이 \((b-a).\mathrm{val}\)의 \texttt{IsPath}이며, \(a=b\)이면 길이 0이다. \(c(t)\)가 이 호의 내부에
있는 조건은 \(0 < (t-a).\mathrm{val} < (b-a).\mathrm{val}\)이다. 따라서 \(a\)에서 \(b\)로 가는 호와
\(b\)에서 \(a\)로 가는 호는 양 끝점을 뺀 내부가 서로소다. 꼭짓점 타입의 유한성은 필요 없다.
\end{lemma}

\begin{proof}
\leanok
표준 사이클 그래프의 걸음을 단사인 꼭짓점 배치로 옮긴다. 원래 순서 방향이 되도록 기본 걸음을 뒤집는다.
시작점을 \(a\)로 회전시킨 뒤 \((b-a).\mathrm{val}\)걸음만 취하면 호가 된다. 전체 길이보다 짧은 사이클의
앞부분은 단순 경로다. 내부의 위치를 모듈러 차의 정수 대표로 판정하면 반대 방향의 두 열린 구간은 겹치지 않는다.
\end{proof}

\begin{lemma}[서로 겹치지 않는 위치 구간의 호]
\label{lem:cycseq_disjoint_arcs}
\uses{lem:cycseq_walks, lem:cyclic}
\lean{Kuratowski.CyclicSequence.disjoint_arc_interiors_of_le,
  Kuratowski.CyclicSequence.disjoint_arc_interiors_wrap,
  Kuratowski.CyclicSequence.not_mem_arc_support_of_not_mem_range,
  Kuratowski.CyclicSequence.not_mem_arc_interior_of_le,
  Kuratowski.CyclicSequence.not_mem_arc_interior_wrap}
\leanok
정수 대표가 \(a.\mathrm{val} \le b.\mathrm{val} \le c.\mathrm{val} \le d.\mathrm{val}\)인 네 위치에 대해,
\(a\)에서 \(b\)로 가는 호와 \(c\)에서 \(d\)로 가는 호의 내부는 서로소다. \(d\)에서 \(a\)로 감기는
호와 \(b\)에서 \(c\)로 가는 호의 내부도 서로소다. 사이클 밖의 점은 어느 호에도 나타나지 않는다.
또 감기지 않는 호의 바깥 위치와 양 끝점은 그 내부에 없고, \(b\)에서 \(a\)로 감기는 호의 내부는
\(a.\mathrm{val} \le t.\mathrm{val} \le b.\mathrm{val}\)인 모든 위치를 피한다.
\end{lemma}

\begin{proof}
\leanok
각 호의 내부를 \texttt{sbtw}의 엄격한 순환 부등식으로 바꾸면 위치 구간들의 부등식이 모순된다.
사이클 밖의 점에 대해서는 호의 모든 점이 순환 수열의 상에 속한다는 사실을 쓴다.
\end{proof}

보조정리~\ref{lem:TK5_of_common}, \ref{lem:TK33_of_alt}에서 이 호들로 세분의 경로 체계를 구성했다.
특히 네 위치의 순환 순서는 서로 반대 방향인 두 열린 호의 \texttt{sbtw} 조건으로 표현한다.
보조정리~\ref{lem:J2}는 꼭짓점 삭제로 생긴 면에서 이 순환 수열과 이웃의 역순을 얻는다.

\section{꼭짓점 나누기}

보조정리~\ref{lem:S}는 수축 \(G/xy\)의 종수 0 회전 체계에서 \(G\)의 종수 0 회전 체계를 만든다. 그 수술을
두 일반 수술로 나눈다. 공통 이웃이 없는 변의 수축을 되돌리는 것(보조정리~\ref{lem:uncontract})과 같은 면의
두 모서리를 변으로 잇는 것(보조정리~\ref{lem:add_edge})이다. 변·고립점·성분 수의 세기는 따로 둔다.
아래에서 다트는 양 끝으로 적는다. Lean에서는 다트 대응과 \texttt{toProd}로 적고 캐스트를 쓰지 않는다.

\begin{lemma}[변을 넣을 때의 세기]
\label{lem:edge_add_counts}
\uses{def:eulerChar}
\lean{Kuratowski.ncard_edgeSet_sup_edge, Kuratowski.ncard_isolated_sup_edge,
  Kuratowski.card_connectedComponent_sup_edge_of_reachable}
\leanok
유한한 꼭짓점 타입 위에서 서로 다르고 인접하지 않은 두 점 사이에 변을 넣으면 변 수는 하나 늘어난다.
두 점이 원래 비고립이면 고립점 수가 같고, 원래 이어져 있으면 성분 수가 같다.
고립점과 성분 수의 등식 자체는 꼭짓점 타입이 무한해도 성립한다.
\end{lemma}

\begin{proof}
\leanok
변 집합에 새 변 하나를 넣고, 이미 비고립인 두 점에만 닿으므로 고립점 집합은 같다.
원래 이어진 두 점 사이의 새 변은 원래 걸음으로 바꿀 수 있어 연결 성분의 동치 관계도 같다.
\end{proof}

\begin{lemma}[공통 이웃이 없는 변의 수축 되돌리기]
\label{lem:uncontract}
\uses{def:contract, def:rotation, def:eulerChar}
\notready
\planned{RotationSystem.exists\_uncontract}
유한하고 같음을 판정할 수 있는 \(V\) 위의 그래프 \(G\)에서 \(x \sim y\)이고 \(x\), \(y\)의 공통 이웃이 없다고 하자.
\(R'\)을 \(G/xy\)의 회전 체계(회전 \(\rho'\))라 하고 \(N\)을 \(G/xy\)에서 \(x\)의 이웃 수라 하자.
\(x\)에서 나가는 \(G/xy\)의 다트 \(d_0\)와 \(0 \le k \le N\)이 있어, \(u_j = (\rho'^{-j} d_0)\)의 끝점
(\(0 \le j < N\))이 \(j < k\)이면 \(G\)에서 \(x\)의 이웃이고 \(k \le j\)이면 \(G\)에서 \(y\)의 이웃이라 하자.
곧 \(x\) 둘레에서 \(G\)의 \(x\)쪽 이웃과 \(y\)쪽 이웃이 각각 한 호를 이룬다. 그러면 \(G\)의 회전 체계 \(R\)이 있어
면 순환 수가 같고(\(O(R) = O(R')\)) 회전이 다음과 같다.
\begin{itemize}
  \item \(x\) 둘레: \(k > 0\)이면 \((x, u_0)\) 다음이 \((x, y)\), \((x, y)\) 다음이 \((x, u_{k-1})\)이고,
        \(0 < j < k\)인 \((x, u_j)\) 다음은 \((x, u_{j-1})\)이다 (\(\rho'\) 그대로). \(k = 0\)이면 \((x, y)\) 하나뿐이다.
  \item \(y\) 둘레: \(k < N\)이면 \((y, u_k)\) 다음이 \((y, x)\), \((y, x)\) 다음이 \((y, u_{N-1})\)이고,
        \(k < j < N\)인 \((y, u_j)\) 다음은 \((y, u_{j-1})\)이다. \(k = N\)이면 \((y, x)\) 하나뿐이다.
  \item \(x\), \(y\)가 아닌 꼭짓점 \(z\) 둘레: \(\rho'\)의 순서에서 \(x\)를, \(z\)가 \(G\)에서 \(y\)의 이웃이면 \(y\)로 바꾼다.
\end{itemize}
다트 대응 \(\{d \in D_G : d \ne (x,y), (y,x)\} \simeq D_{G/xy}\)는 \(y\)를 \(x\)로 바꾸는 사상이며 공개 API로 둔다
(변 수 세기에 다시 쓴다). 세기 \(E(G) = E(G/xy) + 1\), \(I(G) = I(G/xy) - 1\), \(c(G) = c(G/xy) - 1\)과 합치면
\(\chi(R) = \chi(R') - 2\)이므로 \(R'\)이 종수 0이면 \(R\)도 종수 0이다. \(N = 1\)이면 \(k = 0\) 또는 \(k = N\)이다.
\bfcheck{\texttt{[S-분할]}의 첫 단계: 아래 보조정리~\ref{lem:S}의 수술 16,690개 경우(빈 단어 제외)에서
\(G_0\)(아래)의 공통 이웃이 없고 \(G_0/xy = G/xy\)이며, 위 \(R\)이 \(G_0\)의 종수 0 회전 체계이고
\(O(R) = O(R')\)이다. \(x\)쪽 호가 빈 경우 1,658개, \(y\)쪽 호가 빈 경우 6,306개, \(N = 1\)인 경우 1,962개를 포함한다.
\texttt{[S-분할-음성]}: \(y\) 둘레의 \(y\)쪽 호를 뒤집으면 3,030개 경우에서 실패한다.}
\end{lemma}

\begin{proof}
\uses{lem:swap, lem:contract_counts}
\begin{itemize}
  \item \textbf{다트 대응.} \((x,y)\), \((y,x)\)가 아닌 \(G\)의 다트는 끝점의 \(y\)를 \(x\)로 바꿔 \(G/xy\)의 다트로
        보낸다. \(G/xy\)에서 \(y\)는 고립점이고, 공통 이웃이 없으므로 \((x,z)\)와 \((y,z)\)가 같은 다트로 가지 않아
        전단사다. 이 대응은 다트 뒤집기 \(\sigma\)와 가환이다.
  \item \textbf{회전.} \(\tau\)를 대응으로 옮긴 \(\rho'\)에 \((x,y)\), \((y,x)\)를 고정점으로 붙인 순열이라 하자.
        \(u = (x, u_0)\)(\(k > 0\)), \(w = (y, u_k)\)(\(k < N\))로 두면
        \[
          \rho = \tau\,(u\ w)\,(u\ xy)\,(w\ yx) \quad (0 < k < N), \qquad
          \rho = \tau\,(u\ xy) \quad (k = N), \qquad \rho = \tau\,(w\ yx) \quad (k = 0)
        \]
        이다. \(\tau(u) = (y, u_{N-1})\), \(\tau(w) = (x, u_{k-1})\)이므로 \((u\ w)\)는 \(x\)의 \(\tau\)-순환을 \(x\)쪽 호와
        \(y\)쪽 호로 나누고, 뒤의 두 호환은 \(xy\), \(yx\)를 각 호의 끝에 넣는다. 각 꼭짓점의 다트는 한 순환을 이룬다.
  \item \textbf{면.} \((\rho\,(p\ a)(b\ q))\sigma = \rho\sigma\,(\sigma p\ \sigma a)(\sigma b\ \sigma q)\)에서
        \(\varphi = \tau\sigma\,(\bar u\ \bar w)\,(\bar u\ yx)\,(\bar w\ xy)\)이다. \(\tau\sigma\)는 옮긴 \(\varphi'\)에 2-순환
        \((xy\ yx)\)를 더한 것이라 순환 수가 \(O(R') + 1\)이다. \(\bar u\), \(\bar w\)가 \(\varphi'\)의 같은 순환이면 세 호환이
        나누기·합치기·합치기로 \(+1-1-1\), 다른 순환이면 합치기·합치기·나누기로 \(-1-1+1\)이다
        (보조정리~\ref{lem:swap}). 한쪽 호가 비면 호환 하나가 2-순환과 합쳐 \(-1\)이다. 어느 경우든 \(O(R) = O(R')\).
  \item \textbf{세기.} \(d_0\)가 있으므로 \(x\)는 \(G/xy\)와 \(G\)에서 모두 고립점이 아니고, \(y\)는 \(G/xy\)에서만 고립점이다.
        변은 대응에서 \(xy\) 하나가 더 있다. \(G/xy\)의 성분은 \(G\)의 성분에서 \(y\)를 떼어 낸 것과 \(\{y\}\)다.
\end{itemize}
\textbf{아직 형식화할 것.} 다트 대응과 \(\tau\), 세 경우의 회전 식과 한 순환 조건, 면 순열 식과 순환 수 계산,
\texttt{contractEdge}의 변·고립점·성분 수 세기.
\end{proof}

\begin{lemma}[같은 면의 두 모서리를 잇는 변 넣기]
\label{lem:add_edge}
\uses{def:rotation, def:eulerChar}
\notready
\planned{RotationSystem.exists\_addEdge}
유한하고 같음을 판정할 수 있는 \(V\) 위의 그래프 \(G\)와 회전 체계 \(R\)(회전 \(\rho\), 면 순열 \(\varphi\))에서 다트 \(a\),
\(c\)의 시작점 \(u\), \(v\)가 서로 다르고 인접하지 않으며, \(\bar a\)와 \(\bar c\)가 \(\varphi\)의 같은 순환에 있다고 하자.
곧 \(u\)에서 \(a\) 바로 뒤의 모서리와 \(v\)에서 \(c\) 바로 뒤의 모서리가 같은 면에 있다. 그러면 \(G + uv\)
(Lean에서는 \texttt{G} \(\sqcup\) \texttt{edge u v})의 회전 체계 \(R_1\)이 있어 \(a\) 다음이 \((u, v)\), \((u, v)\) 다음이 \(\rho(a)\),
\(c\) 다음이 \((v, u)\), \((v, u)\) 다음이 \(\rho(c)\)이고 나머지 다트의 회전은 \(\rho\)와 같다. 면 순환 수는
\(O(R_1) = O(R) + 1\)이다. 세기 \(E + 1\), 고립점과 성분 수는 그대로(\(u\), \(v\)는 \(a\), \(c\)가 있어 고립점이 아니고
같은 면의 경계를 따라 이어져 있다)와 합치면 \(\chi(R_1) = \chi(R)\)이므로 종수 0이 유지된다.
\bfcheck{\texttt{[S-분할]}의 둘째 단계: 같은 16,690개 경우에서 \(B\) 글자마다 넣는 두 모서리가 같은 면에 있고,
넣은 회전 체계가 종수 0이며 \(O\)가 1 는다. \texttt{[S-분할-음성]}: \(b_1\)을 \(y\) 둘레에서 \(x\) 뒤에 넣으면
4,352개 경우에서 실패한다.}
\end{lemma}

\begin{proof}
\uses{lem:swap, lem:edge_add_counts}
\(G + uv\)의 다트는 \(G\)의 다트와 \((u, v)\), \((v, u)\)다. \(\tau\)를 \(\rho\)에 두 새 다트를 고정점으로 붙인 순열이라
하면 \(\rho_1 = \tau\,(a\ uv)\,(c\ vu)\)이고 \(\varphi_1 = \tau\sigma\,(\bar a\ vu)\,(\bar c\ uv)\)다. \(\tau\sigma\)의 순환 수는
\(O(R) + 1\)이다(2-순환 \((uv\ vu)\)). 첫 호환은 \(\bar a\)의 순환과 2-순환을 합치고(\(-1\)), 합친 순환에는 \(\bar c\)와
\(uv\)가 모두 있으므로 둘째 호환은 나눈다(\(+1\)). 따라서 \(O(R_1) = O(R) + 1\)이다.
\textbf{아직 형식화할 것.} 다트 대응(합집합 다트의 두 다트 제외 전단사), 회전 식과 한 순환 조건, 면 순열 식,
변·고립점·성분 수 세기.
\end{proof}

\begin{lemma}[구간 조건이면 나눠도 평면]
\label{lem:S}
\uses{def:contract, def:IsPlanar, def:rotation}
\notready
\planned{IsPlanar.of\_contractEdge\_of\_interval}
\(x \sim y\)이고 \(R'\)이 \(G/xy\)의 종수 0 회전 체계라 하자. \(x\) 둘레의 \(R'\) 순서로 이웃에 다음 글자를
붙인다.
\begin{itemize}
  \item \(X\): \(G\)에서 \(x\)의 이웃이기만 함
  \item \(Y\): \(y\)의 이웃이기만 함
  \item \(B\): 둘 다
\end{itemize}
이 순환 단어가 \texttt{B? X* B? Y*} 꼴이면 (정규식 표기: \(B\)는 0개 또는 1개, \(X\), \(Y\)는 0개 이상)
\(G\)는 평면이다.

이 꼴의 순환 단어 모임은 회전과 역순에 대해 닫혀 있다. \(b_1 X^p b_2 Y^q\)를 뒤집으면 \(Y^q b_2 X^p b_1\)이고 이를
회전하면 \(b_2 X^p b_1 Y^q\)이다. 그래서 Lean 진술(\texttt{IntervalWord})은 보조정리~\ref{lem:J2}와 같은 역회전 순서로
읽는다. \(N\)을 \(G/xy\)에서 \(x\)의 이웃 수라 할 때, \(N = 0\)이거나, \(x\)에서 나가는 \(G/xy\)의 다트 \(d_0\)와
\(0 \le k \le N\)이 있어 \(u_j = (\rho'^{-j} d_0)\)의 끝점(\(0 \le j < N\))이 다음을 만족한다.
\begin{itemize}
  \item \(j < k\)이면 \(u_j\)는 \(G\)에서 \(x\)의 이웃이다 (\(X\) 또는 \(B\)).
  \item \(k \le j\)이면 \(u_j\)는 \(G\)에서 \(x\)의 이웃이 아니다 (\(Y\)).
  \item \(0 < j\)이고 \(j + 1 < k\)이면 \(u_j\)는 \(G\)에서 \(y\)의 이웃이 아니다 (\(B\)는 \(x\)쪽 호의 양 끝에만).
\end{itemize}
\bfcheck{\texttt{[S]}: 5,307개 경우(연결되지 않은 그래프와 빈 단어 포함). 진술만 본다. 구간 조건이면 \(G\)에
종수 0 회전 체계가 \emph{있는지}를 찾을 뿐, 아래 수술이 만드는 회전 체계를 확인하지는 않는다. 대우에 따라
비평면 그래프에서는 구간 조건이 성립하지 않아야 하는데, 꼭짓점 5개 이하의 비평면 그래프 \(K_5\)는 이미 이
검사에 들어 있다. \texttt{[S-음성]}은 \(K_5\)와 \(K_{3,3}\)의 모든 변 수축과 그 종수 0 회전 체계 76개에서 구간
조건이 실패함을 따로 확인한다. \texttt{[S-구간식]}: 길이 8 이하의 \(X, Y, B\) 단어 9,841개에서 위 역회전 지수의
조건이 정규식 조건과 같고, 정규식 조건은 회전과 역순에 대해 불변이다.}
\end{lemma}

\begin{proof}
\uses{lem:uncontract, lem:add_edge, lem:sum0, lem:small, lem:J2}
\begin{itemize}
  \item \textbf{빈 단어} (\(N = 0\)). \(x\), \(y\)의 \(G\) 이웃은 서로뿐이므로 \(G = (G/xy) \sqcup \{xy\}\)이고 두 지지
        집합은 서로소다 (\(G/xy\)는 \(G\)에서 변 \(xy\)만 뺀 그래프다). \(G/xy\)는 \(R'\)으로 평면이고 변 하나짜리
        그래프는 지지 집합이 두 점이라 평면이다(보조정리~\ref{lem:small}). 보조정리~\ref{lem:sum0}으로 \(G\)는 평면이다.
  \item \textbf{공통 이웃의 자리.} \(N > 0\)이라 하자. 보조정리~\ref{lem:J2}의 이웃 열거로 \(G/xy\)에서 \(x\)의 모든 이웃은
        어떤 \(j < N\)의 \(u_j\)이다. 공통 이웃 \(b\)는 \(x\)의 이웃이므로 \(k \le j\)일 수 없고, \(y\)의 이웃이므로
        \(0 < j < k - 1\)일 수 없다. 따라서 공통 이웃은 \(b_2 = u_0\), \(b_1 = u_{k-1}\) 중에 있다 (\(k \ge 1\)).
        \(k = 1\)이면 \(b_1 = b_2\)다.
  \item \textbf{되돌리기.} \(G_0\)을 \(G\)에서 공통 이웃 \(b\)마다 변 \(yb\)를 뺀 그래프라 하자. \(b\)는 \(G_0\)에서도 \(x\)의
        이웃이므로 \(G_0/xy = G/xy\)이고, \(G_0\)에서 \(x\), \(y\)의 공통 이웃은 없다. \(k \le j\)인 \(u_j\)는 \(G\)에서 \(x\)의
        이웃이 아니고 \(G/xy\)에서 \(x\)의 이웃이므로 \(y\)의 이웃이며, 공통 이웃이 아니라 \(G_0\)에서도 \(y\)의 이웃이다.
        따라서 보조정리~\ref{lem:uncontract}의 가정이 성립하고 \(G_0\)의 종수 0 회전 체계 \(R_0\)을 얻는다.
  \item \textbf{\(b_1\) 넣기} (\(b_1\)이 공통 이웃일 때). \(a = \rho_0^{-1}(y, x)\), \(c = (b_1, x)\)에 보조정리~\ref{lem:add_edge}를
        쓴다. \(\varphi_0(\bar a) = \rho_0(a) = (y, x)\)이고 \(\varphi_0(y, x) = \rho_0(x, y) = (x, b_1) = \bar c\)이므로 같은
        면이다. 결과: \(y\) 둘레에서 \(x\) 바로 앞에 \(b_1\), \(b_1\) 둘레에서 \(x\) 바로 뒤에 \(y\).
  \item \textbf{\(b_2\) 넣기} (\(k \ge 2\)이고 \(b_2\)가 공통 이웃일 때). 앞 단계의 회전 \(\rho_1\)에서 \(a = (y, x)\),
        \(c = \rho_1^{-1}(b_2, x)\)에 보조정리~\ref{lem:add_edge}를 쓴다. \(x\) 둘레는 앞 단계에서 바뀌지 않았으므로
        \(\varphi_1(\bar c) = \rho_1(c) = (b_2, x)\), \(\varphi_1(b_2, x) = \rho_1(x, b_2) = (x, y) = \bar a\)이다. 결과: \(y\) 둘레에서
        \(x\) 바로 뒤에 \(b_2\), \(b_2\) 둘레에서 \(x\) 바로 앞에 \(y\).
  \item 넣은 변을 더하면 \(G\)이고 종수 0이 유지되므로 \(G\)는 평면이다.
\end{itemize}
합성한 회전은 다음과 같다(\(\rho'\) 순서로 \(x\) 둘레 단어를 \(b_1 X_1 \cdots X_p\, b_2\, Y_1 \cdots Y_q\)로 읽음).
\begin{itemize}
  \item \(x\) 둘레: \((b_1, X_1, \ldots, X_p, b_2, y)\)
  \item \(y\) 둘레: \((b_2, Y_1, \ldots, Y_q, b_1, x)\)
  \item \(b_1\) 둘레: \(R'\)에서 \(x\)가 있던 자리에 \(x, y\)를 이 순서로 넣는다.
  \item \(b_2\) 둘레: \(x\)가 있던 자리에 \(y, x\)를 이 순서로 넣는다.
  \item \(Y\) 글자인 이웃 둘레: \(x\)를 \(y\)로 바꾼다.
  \item 그 밖의 꼭짓점 둘레: \(R'\)과 같다.
\end{itemize}
면 수로 보면 \(x\)를 지나지 않는 \(R'\)의 면은 그대로 남고, \(B\) 글자 \(b\)마다 삼각형 면이 하나씩 생긴다
(예: \(b_1\)에서 \((y, x) \mapsto (x, b_1) \mapsto (b_1, y) \mapsto (y, x)\)). 곧 \(O = O' + |B|\),
\(E = E' + 1 + |B|\), \(I = I' - 1\), \(c = c' - 1\)이고 \(\chi = \chi' - 2 = 2c\)다.
\bfcheck{\texttt{[S-수술]}: 16,880개 경우(빈 단어 190개 포함)에서 위 합성 회전이 \(G\)의 회전 체계이고
종수 0이며 \(O = O' + |B|\)(빈 단어이면 \(O' + 1\))이다. \texttt{[S-수술-음성]}: \(b_1\) 둘레의 순서를
\(y, x\)로 바꾸면 5,454개 경우에서, \(y\) 둘레에서 \(b_1\)과 \(b_2\)를 바꾸면 5,908개 경우에서 반례가 나온다.
\texttt{[S-분할]}: 빈 단어를 뺀 16,690개 경우에서 되돌리기와 변 넣기의 합성이 위 회전과 정확히 같다.
\(B\)가 하나뿐이고 \(b_1 = b_2\)인 경우 2,982개를 포함한다.}

\textbf{아직 형식화할 것.} 보조정리~\ref{lem:uncontract}, \ref{lem:add_edge}, 빈 단어의 서로소 합,
공통 이웃의 자리, \(G_0\)과 그래프 등식 \(G_0/xy = G/xy\), 변을 넣은 뒤 \(G\)와의 그래프 등식. 주 정리에서
보조정리~\ref{lem:cyclic}의 (a)를 이 구간 조건으로 바꾸는 다리는 후속으로 형식화한다.
\end{proof}

\section{경우 나누기}

\begin{lemma}[순환 단어]
\label{lem:cyclic}
\lean{Kuratowski.CyclicWord.Letter, Kuratowski.CyclicWord.xPositions,
  Kuratowski.CyclicWord.yPositions, Kuratowski.CyclicWord.bothPositions,
  Kuratowski.CyclicWord.Consecutive, Kuratowski.CyclicWord.Clustered,
  Kuratowski.CyclicWord.Alternating, Kuratowski.CyclicWord.classify,
  Kuratowski.CyclicWord.sbtw_iff_val, Kuratowski.CyclicWord.mem_walkInterior_arc_iff}
\leanok
순환 단어의 글자는 \(X\), \(Y\), \(B\), \(-\)이고, \(X/B\)인 글자와 \(Y/B\)인 글자가 각각 2개 이상이라
하자. 그러면 다음 중 하나가 성립한다.
\begin{enumerate}
  \item[(a)] \(Y/B\)인 글자가 모두 연속한 두 \(X/B\) 글자 사이의 닫힌 구간에 있다.
  \item[(b)] \(B\)가 3개 이상이다.
  \item[(c)] \(Y/B\)인 \(u\), \(w\)가 있어, \(u\)와 \(w\)를 뺀 두 열린 호 각각에 \(X/B\) 글자가 있다.
\end{enumerate}
(c)의 호는 끝점 \(u\), \(w\)를 뺀 열린 호다. 닫힌 호로 읽으면 \(u\)나 \(w\)가 \(B\)일 때 끝점만으로 (c)가
성립해 버려, \(K_{3,3}\)을 만드는 데 필요한 내용이 사라진다.
\bfcheck{\texttt{[cyclic]}: 길이 2 이상 9 이하 단어 331,476개. \texttt{[cyclic}\(\equiv\)\texttt{S]}: 29,330개}
\end{lemma}

\begin{proof}
\leanok
\(X/B\) 글자들은 단어를 구간들로 나눈다. 연속한 두 \(X/B\) 글자 \(a\), \(a'\) 사이의 열린 구간을 \((a, a')\),
닫힌 구간을 \([a, a']\)라 하자. (a)와 (b)가 모두 성립하지 않는다고 하고 (c)를 보인다.
\begin{itemize}
  \item \(Y\) 글자 \(u\)가 있으면 \(u\)는 어떤 \((a, a')\) 안에 있다. (a)가 거짓이므로 \([a, a']\) 밖에 \(Y/B\)
        글자 \(w\)가 있다. \(u\)는 \(X/B\)가 아니고 \(w\)는 \([a, a']\) 밖에 있으므로, \(u\)와 \(w\)를 뺀 두 열린
        호는 각각 \(a\)와 \(a'\)를 지난다.
  \item \(Y\) 글자가 없으면 \(Y/B\) 글자는 모두 \(B\)이고, (b)가 거짓이므로 정확히 2개다. 이 둘을 \(u\), \(w\)라
        하자. (a)가 거짓이므로 \(u\)와 \(w\)는 연속한 \(X/B\) 글자가 아니고, 따라서 둘 사이의 두 열린 호에 모두
        \(X/B\) 글자가 있다.
\end{itemize}
Lean에서는 양의 길이 \(m\)의 위치를 \texttt{ZMod m}으로 나타내고 Mathlib의 순환 순서
\texttt{btw}, \texttt{sbtw}로 닫힌 호와 열린 호를 표현했다. \(X/B\), \(Y/B\) 위치를 두 유한 집합으로
두면 그 교집합이 \(B\) 위치다. \texttt{classify}는 네 글자 단어에 이 분류를 적용한다.
길이 2도 허용하며, 두 글자 집합이 각각 두 점 이상이라는 가정은 그대로다.
\texttt{sbtw}가 나타내는 열린 호는 앞의 실제 경로 내부와 같다는 것도 증명했다.
\end{proof}

이웃만 남긴 단어에서 (a)는 보조정리~\ref{lem:S}의 정규식 구간 조건과 같다. 이 정규식과의 동치
자체는 아직 Lean으로 형식화하지 않았고, 위 \texttt{[cyclic}\(\equiv\)\texttt{S]} 검사로만 확인했다.
